\section{Related work}\label{sec:related}

\subsection{Deep multivariate forecasting}

A line of work argues that carefully regularised, architecturally simple models match or beat transformers on multivariate series. N-BEATS~\cite{nbeats} and N-HiTS~\cite{nhits} stack residual MLP blocks with backcast/forecast decomposition; DLinear~\cite{dlinear} reduces forecasting to two linear maps on a trend/remainder split and questions whether attention helps at all; TSMixer~\cite{tsmixer} alternates time- and feature-mixing MLPs. Against these, Autoformer~\cite{autoformer} introduces series decomposition with FFT auto-correlation, PatchTST~\cite{patchtst} treats a series as patches with channel independence and reversible instance normalisation, TimesNet~\cite{timesnet} folds FFT-selected periods into 2-D convolutions, and iTransformer~\cite{itransformer} inverts the attention axis so that variates rather than time steps are tokens. These models are evaluated almost exclusively on long, densely sampled, quasi-stationary benchmarks (traffic, electricity, weather) with thousands of time steps. Our panels have 37 and 99 weeks and 57--74\% zeros, and our results reproduce the simplicity finding in that regime: N-BEATS and N-HiTS are among the strongest baselines, while TimesNet, Autoformer and TSMixer fall below a naive forecast on at least one panel.

\subsection{Spatio-temporal graph neural networks}

Where regions interact, graph structure is the natural inductive bias. AGCRN~\cite{agcrn} learns node-adaptive parameters and an adjacency from node embeddings inside a GRU; MTGNN~\cite{mtgnn} learns a sparse directed graph jointly with dilated temporal convolutions and mix-hop propagation; STID~\cite{stid} shows much of the benefit can be recovered with spatial and temporal identity embeddings plus an MLP; STAEformer~\cite{staeformer} adds an adaptive spatio-temporal embedding to a vanilla transformer. A recurring question in this literature is whether the learned graph encodes anything beyond a smoothed distance kernel. Our gating analysis speaks to this directly: the learned and attention channels together receive only 24\% of the mixing weight on WA2014, against 47\% for the fixed geographic kernel and 28\% for the genomic graph, and the genomic channel is shown to be \emph{negatively} correlated with distance, so it is the externally supplied prior rather than the learned structure that adds the long-range information.

\subsection{Epidemic-specific graph models}

Epidemiological GNNs add domain structure to the above. Cola-GNN~\cite{colagnn} combines cross-location attention with a geographic adjacency through a learned gate, an idea our fusion generalises to four channels. EpiGNN~\cite{epignn} encodes local and global transmission risk and a region-aware graph learner. MepoGNN~\cite{mepognn} is the closest prior work in spirit: it couples a spatio-temporal network to a metapopulation SIR model with a learned mobility graph, so that the network predicts epidemiological parameters rather than counts. Our decoder differs in the mechanism it couples to. SIR compartments require a population denominator and a recovery rate, which are weakly identified for Ebola at health-zone resolution; empirically MepoGNN is the weakest graph baseline on both our panels. The renewal equation instead requires only the generation-interval distribution, which for Ebola is measured~\cite{who2014}, and it is the estimating equation already used in outbreak response.

\subsection{Renewal-equation estimation}

The renewal formulation~\cite{fraser2007} and its practical estimator EpiEstim~\cite{cori2013} are standard for estimating a time-varying reproduction number, and metapopulation extensions couple regions through mobility~\cite{riley2007}. This literature is estimation-oriented: $R_t$ is inferred retrospectively from observed incidence under a fixed generation interval, usually without a learned spatial operator and without multi-step forecasting. We invert the usual arrangement by making the renewal equation a differentiable \emph{decoder} whose parameters are predicted by a GNN, which keeps the interpretability of $R_t$ while producing autoregressive multi-week forecasts. The price, as our theory and experiments both show, is that the decoder inherits the renewal equation's momentum: it contracts when $\bar R<1$ and amplifies when $\bar R>1$.

\subsection{Genomic epidemiology}

Phylodynamic analysis of the 2013--16 West African epidemic~\cite{dudas2017} reconstructed the spread of 1{,}610 Makona genomes and identified cross-border introductions and the role of mobility and administrative boundaries; genomic surveillance was integrated into the Nord Kivu response in near-real time with 792 genomes analysed as they were generated~\cite{kinganda2021}. These studies use genomes to explain transmission retrospectively. The enabling capability is portable real-time sequencing, demonstrated during the West African epidemic itself~\cite{quick2016} and since generalised into routine outbreak genomic surveillance~\cite{grubaugh2019}, with open platforms such as Nextstrain making lineage information available while an outbreak is still running~\cite{hadfield2018}. That these data now arrive in near-real time is precisely what makes a forecasting prior conceivable; what remains missing is a model that consumes them prospectively. To our knowledge genomes have not previously been used as a \emph{forecasting} prior, and doing so raises a leakage problem absent from retrospective work: a phylogeny estimated from the full data set encodes the future. We therefore rebuild the linkage graph at every forecast origin from genomes sampled strictly before it, and report how sequencing breadth across regions, rather than depth, determines how much of that graph is usable.

\subsection{Evaluation practice}

Forecasting comparisons are sensitive to protocol. The predictive median is the Bayes act under absolute loss~\cite{gneiting2011}, a point our ablations confirm is worth 12--13\% of validation MAE on these counts, and proper scoring of the predictive distribution is standard in collaborative epidemic forecasting, where the weighted interval score has become the common currency~\cite{bracher2021} and large multi-team evaluations have shown how much conclusions depend on the scoring protocol~\cite{cramer2022}. Pairwise significance in this setting is conventionally assessed with the Diebold--Mariano test~\cite{diebold1995}, which we use alongside a rank-based correction. Much of the spatio-temporal GNN literature reports a single chronological split, which on an epidemic panel can place the entire test period in a near-zero tail; our rolling-origin protocol with phase-labelled blocks is designed to avoid that, and our measured validation-to-test rank correlation of 0.72 on WA2014 quantifies how badly a single split can mislead.
